\section{Discussion}
\label{sec:discussion}

\subsection{Tractable Hamiltonian Classes}
\label{subsec:tractable}

The eigenbasis requirement
(Section~\ref{subsec:eigenbasis}) restricts stereographic block
encoding to structured Hamiltonians.  The three classes identified
in that section---computational-basis diagonal operators,
Hamiltonians with known diagonalization, and QPE-assisted
encoding---cover a range of physically relevant systems.

Diagonal operators arise naturally in quantum chemistry
(diagonal Coulomb integrals), combinatorial optimization
(QUBO/Ising cost functions), and classical simulation.  For
these, stereographic block encoding provides a concrete advantage:
the encoding depends only on individual eigenvalues, not on the
global norm, and the single-ancilla construction is simpler than
standard block encoding.

Hamiltonians with known diagonalization include integrable spin
chains solvable by the Bethe ansatz~\cite{Bethe1931}, free-fermion
models diagonalizable by the Jordan--Wigner
transformation~\cite{JordanWigner1928}, and small systems where
exact diagonalization is tractable.  The Heisenberg model
(Section~\ref{subsec:heisenberg}) is a representative example:
the change of basis is a Givens rotation with $\mathcal{O}(1)$
gate complexity.

% TODO: Expand discussion of integrable models.
% Which Bethe-ansatz models have efficient circuit implementations
% of the diagonalization?

\subsection{Connection to Quantum Phase Estimation}
\label{subsec:qpe-connection}

Quantum phase estimation (QPE) can, in principle, extract the
eigenvalues of a Hamiltonian into a register, enabling conditioned
operations on each eigenspace.  This suggests a hybrid approach:
use QPE to resolve eigenspaces, then apply stereographic encoding
conditioned on the QPE register.  However, this reintroduces the
resources that stereographic encoding aims to avoid: QPE requires
$\mathcal{O}(\log(1/\epsilon))$ ancilla qubits and
$\mathcal{O}(\alpha/\epsilon)$ queries to a block encoding
of~$H/\alpha$.  The advantage of stereographic encoding over
standard QSVT in this hybrid setting is unclear and deserves
further investigation.

% TODO: Formalize the hybrid QPE + stereo-QSP protocol.
% Under what conditions does it outperform standard QSVT?

\subsection{Generalized Signal Operators}
\label{subsec:generalized}

Replacing the phase rotations $\exp[i\phi\PauliZ]$ in the QSP
product with general signal-type operators $O_\alpha$
parameterized by $\alpha \in \hat{\C}$ yields a generalized
normalization condition
\begin{equation}
	|P|^2 + |Q|^2
	= (1{+}r^2)^d \textstyle\prod_j(1{+}|\alpha_j|^2),
	\label{eq:generalized-norm}
\end{equation}
whose forward direction follows from the same inductive argument
as the unnormalized formulation
(Section~\ref{subsec:review-function-class}), but whose
converse---decomposing arbitrary polynomial pairs into such a
product---remains open.  This generalization could enlarge the
class of achievable functions and relax the parity constraint.

\subsection{Comparison with Related Work}
\label{subsec:related}

Stereographic block encoding is complementary to several recent
extensions of the QSP framework.
GQSP~\cite{MotlaghWiebe2023} lifts the parity and realness
constraints and provides faster angle-finding, but retains
bounded output.  Hamiltonian block encoding~\cite{KangSu2025}
uses at most two ancilla qubits without explicit normalization,
but still produces bounded polynomial transformations.
Randomized QSVT~\cite{WangZhangHazraEtAl2025} avoids block
encodings through importance sampling, addressing a different
cost bottleneck.  Stereographic encoding is the only approach
that produces genuinely unbounded output, at the cost of the
eigenbasis requirement and increased sampling overhead.

\subsection{Open Problems}
\label{subsec:open}

Several questions remain open:
\begin{enumerate}
	\item Can the eigenbasis requirement be relaxed for specific
	      Hamiltonian classes, perhaps by exploiting partial
	      spectral information (e.g., knowledge of the ground
	      state but not the full eigenbasis)?
	\item What is the optimal measurement strategy for
	      extracting $f(\lambda_j)$ from a superposition of
	      eigenspaces?  Can adaptive measurement schemes improve
	      the $\mathcal{O}(r^6/\epsilon^2)$ scaling?
	\item Does the generalized normalization
	      condition~\eqref{eq:generalized-norm} admit a
	      completeness theorem analogous to the standard QSP
	      theorem?
	\item Can stereographic encoding be combined with
	      tensor-network techniques for Hamiltonians with
	      known matrix product operator (MPO) structure?
\end{enumerate}
